\documentclass[a4paper]{article}
% Español (México) con XeLaTeX: fontspec para fuentes Unicode, polyglossia
% para la división silábica y los nombres del documento en la variante mexicana.
\usepackage{fontspec}
\usepackage{polyglossia}
\setdefaultlanguage[variant=mexican]{spanish}
% Los nombres chinos van como «pinyin (original)»: xeCJK da a los caracteres
% Han su propia fuente; sin ella XeLaTeX los omite con un aviso.
% FandolSong no cubre algunos caracteres de nombres (U+73FA); WenQuanYi Zen Hei sí.
\usepackage[AutoFallBack=true]{xeCJK}
\setCJKmainfont{FandolSong-Regular.otf}
\setCJKfallbackfamilyfont{\CJKrmdefault}{WenQuanYi Zen Hei}
% polyglossia no traduce los nombres de listings.
\AtBeginDocument{\providecommand{\lstlistingname}{}\renewcommand{\lstlistingname}{Listado}%
  \providecommand{\lstlistlistingname}{}\renewcommand{\lstlistlistingname}{Índice de listados}}
\usepackage{amsmath, amssymb}
\usepackage{graphicx}
\usepackage[margin=2.5cm]{geometry}
\usepackage[most]{tcolorbox}
\usepackage{etoolbox}
\usepackage{listings}
\usepackage{hyperref}
\usepackage{booktabs}
\usepackage{subcaption}
\usepackage{float}

\newtcolorbox{knowledgebox}[1]{enhanced, colback=blue!5!white, colframe=blue!75!black, colbacktitle=blue!75!black, coltitle=white, fonttitle=\bfseries, title={#1}, attach boxed title to top left={yshift=-2mm, xshift=2mm}, boxrule=1pt, sharp corners}
\newtcolorbox{importantbox}[1]{enhanced, colback=yellow!10!white, colframe=yellow!80!black, colbacktitle=yellow!80!black, coltitle=black, fonttitle=\bfseries, title={#1}, sharp corners}
\newtcolorbox{warningbox}[1]{enhanced, colback=red!5!white, colframe=red!75!black, colbacktitle=red!75!black, coltitle=white, fonttitle=\bfseries, title={#1}, sharp corners}

\lstset{language=Python, basicstyle=\ttfamily\small, keywordstyle=\color{blue}, stringstyle=\color{red!60!black}, commentstyle=\color{green!60!black}, breaklines=true, frame=single, numbers=left, numberstyle=\tiny\color{gray}, captionpos=b, extendedchars=true}

\newcommand{\notetitle}{{[}LLM Agents F24{]} Compound AI Systems \& the DSPy Framework — Omar Khattab}
\newcommand{\noteauthors}{Elaboradas a partir de materiales públicos del curso}
\newcommand{\notedate}{7 de octubre de 2024}
\newcommand{\videochannel}{Berkeley RDI}
\newcommand{\videopublishdate}{7 de octubre de 2024}
\newcommand{\videoduration}{aproximadamente 90 minutos}
\newcommand{\videourl}{https://www.youtube.com/live/JEMYuzrKLUw}
\newcommand{\videocoverpath}{cover.jpg}
\newcommand{\repourl}{https://github.com/hqhq1025/ai-course-notes}


% Footer with repo info
\usepackage{fancyhdr}
\pagestyle{fancy}
\fancyhf{}
\fancyfoot[C]{\thepage}
\fancyfoot[R]{\footnotesize\color{gray}\href{\repourl}{AI Course Notes}}
\renewcommand{\headrulewidth}{0pt}
\renewcommand{\footrulewidth}{0pt}

\begin{document}
\begin{titlepage}\centering
{\Large Notas del curso\par}\vspace{1.2cm}
{\huge\bfseries \notetitle\par}\vspace{0.8cm}
{\large \noteauthors\par}\vspace{0.3cm}
{\large \notedate\par}\vspace{1.2cm}
\ifdefempty{\videocoverpath}{{\small Favor de indicar la ruta de la imagen de portada.\par}}{\includegraphics[width=0.82\textwidth,height=0.45\textheight,keepaspectratio]{\videocoverpath}\par}
\vfill
\begin{tcolorbox}[width=0.9\textwidth, colback=black!2!white, colframe=black!60, sharp corners]
\textbf{Autor o canal del video}: \videochannel\par
\textbf{Fecha de publicación}: \videopublishdate\par
\textbf{Duración del video}: \videoduration\par
\ifdefempty{\videourl}{}{\textbf{Enlace del video}: \href{\videourl}{\nolinkurl{\videourl}}\par}\par
\textbf{Página del proyecto}: \href{\repourl}{\nolinkurl{\repourl}}\par
\end{tcolorbox}
\end{titlepage}

\tableofcontents\newpage


\section{Introducción: de un solo modelo a sistemas de IA compuestos}

Omar Khattab es científico investigador en Databricks y creador del framework DSPy. La idea central de esta clase: \textbf{una sola llamada a un LLM no es suficiente}; necesitamos combinar múltiples llamadas a LLM, recuperación de información y herramientas en \textbf{sistemas de IA compuestos} (Compound AI Systems), y construirlos mediante \textbf{programación} en lugar de prompting manual.

\begin{importantbox}{De la demo a la producción}
Construir una demo de IA impresionante nunca había sido tan fácil (como ChatGPT), pero construir un sistema de producción \textbf{confiable} sigue siendo difícil. La razón central: la capacidad de una sola llamada a un LLM es limitada y poco confiable.
\end{importantbox}

\subsection{Resumen de la sección}

Los sistemas de IA compuestos, al combinar múltiples componentes, superan el límite de capacidad de una sola llamada a un LLM.
\section{El paradigma de Compound AI Systems}

\subsection{¿Qué es Compound AI Systems?}

\begin{knowledgebox}{Sistema de IA compuesto vs. modelo único}
Un sistema de IA compuesto está formado por múltiples componentes que interactúan entre sí: llamadas a LLM, recuperadores, ejecutores de código, verificadores, etc. Cada componente tiene su propia firma (tipos de entrada/salida) y se combinan mediante lógica de programa. Esto es similar a la «composición de funciones» de la ingeniería de software tradicional.
\end{knowledgebox}

\begin{itemize}
    \item \textbf{RAG} es el sistema compuesto más simple: recuperación + síntesis mediante LLM
    \item \textbf{ReAct} es un sistema compuesto cíclico: iteración de razonamiento y llamadas a herramientas
    \item \textbf{Multi-hop QA}: encadenamiento de múltiples rondas de recuperación y razonamiento
    \item Sistemas más complejos: pipeline de múltiples etapas, que incluye pasos de verificación, reflexión, etc.
\end{itemize}

\subsection{Los problemas del prompting manual}

\begin{warningbox}{Prompting no es sostenible}
Escribir un prompt a mano parece simple, pero en sistemas compuestos provoca:
\begin{itemize}
    \item \textbf{Fragilidad}: cambiar de modelo o modificar el pipeline exige reescribir todos los prompts
    \item \textbf{No composable}: es difícil combinar de forma modular múltiples prompts
    \item \textbf{Difícil de optimizar}: no es posible buscar de forma sistemática el prompt óptimo
    \item \textbf{No portable}: un buen prompt para un modelo puede no funcionar en otro
\end{itemize}
\end{warningbox}

\subsection{Resumen de la sección}

El prompting manual no es sostenible en sistemas compuestos; se necesita un paradigma de programación más abstracto.

\section{DSPy: programación declarativa de IA}

\subsection{Idea central}

\begin{importantbox}{La idea central de DSPy}
DSPy permite que los desarrolladores definan la estructura y la lógica de un sistema de IA mediante un \textbf{programa de Python}, y el marco de trabajo \textbf{compila} automáticamente el prompt y los ejemplos few-shot óptimos. Los desarrolladores solo necesitan declarar «qué hacer» (what); no necesitan especificar manualmente «cómo hacerlo» (how, es decir, el texto concreto del prompt).
\end{importantbox}

Abstracciones clave:
\begin{itemize}
    \item \textbf{Signature}: define los tipos de entrada y salida de un módulo (por ejemplo, «question -> answer»)
    \item \textbf{Module}: componente de IA componible (como ChainOfThought, ReAct)
    \item \textbf{Optimizer}: busca automáticamente el prompt, los ejemplos few-shot y la configuración de módulos óptimos
    \item \textbf{Metric}: define el criterio de evaluación de la calidad del sistema
\end{itemize}

\subsection{Optimizadores (Optimizers)}

DSPy ofrece varios optimizadores para mejorar el sistema de manera automática:

\begin{itemize}
    \item \textbf{BootstrapFewShot}: selecciona automáticamente los mejores ejemplos few-shot a partir de los datos de entrenamiento
    \item \textbf{COPRO}: optimiza automáticamente el texto de instrucción dentro del prompt
    \item \textbf{MIPRO}: optimiza simultáneamente las instrucciones y los ejemplos
    \item \textbf{BootstrapFinetune}: genera datos de entrenamiento y aplica fine-tuning al modelo subyacente
\end{itemize}

\subsection{Compilación frente a escritura manual}

\begin{knowledgebox}{Analogía con la programación tradicional}
DSPy es a prompt engineering lo que un compilador es al lenguaje ensamblador. El desarrollador escribe el programa con abstracciones de alto nivel, y el compilador (optimizer) genera automáticamente la implementación subyacente (el prompt). Esto no solo mejora la eficiencia, sino que también dota al sistema de portabilidad: para cambiar de modelo basta con volver a compilar.
\end{knowledgebox}

\subsection{Resumen de la sección}

Mediante la programación declarativa y la optimización automática, DSPy eleva la construcción de sistemas de IA compuestos, de la prompt engineering manual, a una ingeniería de software sistemática.

\section{Reglas empíricas de los Compound AI Systems}

\subsection{Comparación de los patrones principales}

\begin{table}[H]
\centering
\caption{Patrones de construcción de los Compound AI Systems}
\begin{tabular}{p{0.2\textwidth}p{0.26\textwidth}p{0.22\textwidth}p{0.24\textwidth}}
\toprule
Patrón & Composición & Ventaja & Reto de ingeniería \\
\midrule
RAG & Recuperación + 1 llamada a LLM & Implementación sencilla & Carece de capacidad de razonamiento encadenado \\
ReAct & Ciclo de razonamiento → herramienta → observación & Adecuado para tareas de varios pasos & Requiere una lógica de control minuciosa \\
DSPy & Estructura declarativa + optimizador & Programable y transferible & Espacio de búsqueda grande, el ajuste consume mucho tiempo \\
Autonomous Agent & Colaboración de varios Agent + MCP & Flujos de trabajo paralelos automatizados & Requiere monitoreo fuerte y barreras de seguridad \\
\bottomrule
\end{tabular}
\end{table}

\begin{knowledgebox}{Reglas empíricas}
El diseño de los sistemas de IA compuestos sigue tres principios: 1) los componentes necesitan una firma clara para poder componerse; 2) la estructura debe verificarse con metric / log; 3) el proceso de optimización debe ser reproducible, o de lo contrario el sistema es propenso al drift.
\end{knowledgebox}

\subsection{Firma y composición}

Signature define la forma de entrada/salida de un módulo; DSPy usa la verificación de signature para evitar combinaciones no válidas. Cada vez que se amplía el pipeline, se debe insertar un verificador de tipos y generar automáticamente pruebas de interacción, para evitar que los módulos posteriores reciban valores anómalos.

\subsection{Reutilización de patrones}

Los patrones comunes de los sistemas compuestos (como retriever → reasoner, planner → executor) pueden encapsularse como reusable module. Una vez que la firma es consistente, pueden reutilizarse de manera recurrente en distintos pipeline, lo que reduce considerablemente el costo de copiar prompts a mano.

\begin{knowledgebox}{La reutilización exige firmas consistentes}
Ni el optimizador más potente puede proteger las combinaciones que no coinciden. Si un módulo ``state → action'' se inserta por error en un pipeline ``question → answer'', se generarán continuamente invalid token en la Evidence Matrix. El verificador de firmas debe lanzar un error ya en la etapa de compile.
\end{knowledgebox}

\begin{warningbox}{No trates el pipeline como una caja negra}
En un pipeline que incluye recuperación, razonamiento, herramientas y validación, un error en cualquier paso puede propagarse, así que colocar guard rails (aserciones, verificación de tipos, fallback) en cada límite de módulo es más importante que optimizar el prompt.
\end{warningbox}

\subsection{Resumen de la sección}

La confiabilidad de los Compound AI Systems proviene de interfaces claras, la verificación de que los componentes puedan combinarse y un proceso de optimización observable, no solo de depender de un LLM más potente.
\section{Práctica de compilación y depuración de DSPy}

\subsection{Registros y línea de tiempo}

El optimizer de DSPy genera múltiples combinaciones de prompt; se recomienda dividir el flujo de registros en tres etapas:

\begin{enumerate}
    \item \textbf{Etapa de entrada}: registra la signature, los ejemplos few-shot y la metric esperada
    \item \textbf{Etapa de compilación}: guarda los prompt candidatos generados por el optimizer y sus puntuaciones
    \item \textbf{Etapa de ejecución}: compara el desempeño de cada candidato con datos reales y registra el trace y la response
\end{enumerate}

\subsection{Consejos de depuración}

Para depurar el compilador de DSPy pueden usarse los siguientes métodos:

\begin{itemize}
    \item Formatear cada salida del optimizer como YAML legible para humanos (human-readable) y anotar el efecto
    \item Insertar un flujo fallback: si algún módulo produce una salida anómala, retroceder automáticamente al último pipeline exitoso
    \item Conservar en la Evidence Matrix la correspondencia prompt → response → metric, para facilitar un rastreo rápido
\end{itemize}

\begin{knowledgebox}{La Evidence Matrix es el registro de depuración}
La Evidence Matrix muestra los prompt candidatos, las respuestas generadas y las puntuaciones; cuantas más versiones históricas se registren, más fácil será localizar en qué paso se originó el sesgo. Incorporar una captura de esta matriz a la nota (como [cover.jpg](/home/v-haoqiwang/ai-course-notes/talks/berkeley-llm-agents/f24/lecture08/cover.jpg)), aunque no sea un fotograma en ejecución, también transmite la idea de que «la compilación es observable».
\end{knowledgebox}

\subsection{Evidence Frame}

\begin{figure}[H]
\centering
\includegraphics[width=0.85\textwidth]{cover.jpg}
\caption{Ejemplo de Evidence Matrix: tabla de registro de prompt → resultado → metric}
\end{figure}

\subsection{Resumen de la sección}

La compilación de DSPy necesita registrar cada candidato durante el proceso de optimización de manera que pueda reproducirse; la Evidence Matrix es la fuente primaria para la depuración y la reversión.

\section{Monitoreo, observabilidad y gobernanza}

\subsection{Señales de monitoreo de despliegue}

\begin{table}[H]
\centering
\caption{Lista de verificación de monitoreo de Deployment}
\begin{tabular}{p{0.25\textwidth}p{0.28\textwidth}p{0.25\textwidth}}
\toprule
Señal & Significado & Acción activada \\
\midrule
Latency increase & Degradación del desempeño del modelo o la herramienta & Revertir al last-known-good prompt y agregar una optimizer sample \\
Tool error & Respuesta anómala de la herramienta & Suspender el módulo correspondiente, emitir una alerta y hacer una verificación secundaria \\
Metric drift & Caída en la puntuación & Volver a recopilar los datos de entrenamiento e iniciar el auto-optimizer \\
User feedback spike & Aumento súbito en la tasa de rechazo de respuesta o de malentendidos & Incorporar una revisión humana y actualizar el evidence trace \\
\bottomrule
\end{tabular}
\end{table}

Estas señales normalmente son ingested al observability stack. Latency y Tool error hacen que el alerting se active de inmediato; Metric drift, en cambio, necesita schedule una ejecución comparativa (por ejemplo, un nightly run). User feedback spike normalmente usa QA note + Slack channel (como `\#agent-monitors`) para hacer triage manual.

\begin{warningbox}{Ciclo cerrado de gobernanza}
En cada despliegue importante hay que ejecutar el three-step workflow: define → monitor → annotate. El paso annotate es fundamental, pues asegura que el optimizer posterior siga recordando por qué se configuró así en un inicio.
\end{warningbox}

\subsection{Action Timeline}

En el README de cada release hay que actualizar la tabla de Action Timeline, donde se enumeran el nombre del módulo, la condición que lo activa, el metric threshold y el responsible owner. De esta manera, cuando surge un problema es posible ubicar rápidamente la cadena de responsabilidad.

Action Timeline también puede vincularse con la firma del colaborador: cada ajuste del pipeline debe llevar ``who changed what'', lo que permite deslindar responsabilidades rápidamente tras un incidente. Para equipos multiinquilino, se recomienda guardar también cada timeline entry en el central Incident Log, y enlazarlo mediante el `Evidence Matrix ID`.

\subsection{Resumen de la sección}

El monitoreo y la gobernanza son la tercera capa de protección del DSPy pipeline: es necesario convertir señales como latency, metric y feedback en respuestas ejecutables, en lugar de dejarlas sin control.

\section{Flujo de trabajo y simulacros de incidentes}

\subsection{Simulacro antes del despliegue}

Antes de cada deployment, debe simularse una ronda de tarea representativa: desde dataset, signature, optimizer hasta los resultados de evidence. El objetivo del simulacro es asegurar que el equipo de DevOps y el prompt obtenido por el optimizer se mantengan alineados.

El proceso del simulacro conviene realizarlo en un entorno de staging: se sube el evidence snapshot al dashboard, para que el equipo de QA verifique de antemano el metric de cada module. Solo cuando el evidence run coincide con el production run (con un error < 5 %) se permite pasar a producción.

\subsection{Proceso de respuesta a incidentes}

El proceso de respuesta a incidentes consta de tres pasos: 1) trace → 2) rollback → 3) annotate. La etapa de trace debe localizar con rapidez qué module produjo el drift; la etapa de rollback debe fijar el last-known-good prompt; la etapa de annotate debe registrar el evento en el governance log.

Al reproducir el trace debe adjuntarse el timeline (Action Timeline / Evidence Matrix) y compararse las marcas de tiempo del metric drift. Esto ayuda al equipo a entender si el drift proviene de la entrada del optimizer o de la salida del tool.

\begin{importantbox}{Checklist del simulacro de incidentes}
\begin{itemize}
    \item Confirmar la condición que dispara la alarm (latency/metric/feedback drift)
    \item Pausar el optimizer y conservar el evidence snapshot
    \item Registrar los recovery steps en el Action Timeline, incluyendo owner y timestamp
\end{itemize}
\end{importantbox}

\subsection{Resumen de la sección}

El flujo de trabajo y los simulacros de incidentes llevan el DSPy pipeline del nivel de research al de producción, y mediante el checklist y el Action Timeline permiten al equipo recuperarse con rapidez cuando ocurre un incidente.
\section{Métricas de conversión en producto y ROI}

\subsection{Lista de métricas}

Después de poner en producción el sistema DSPy es necesario recopilar de manera simultánea las siguientes métricas: costo de LLM por consulta, tasa de éxito de invocación de herramientas, tiempo de ejecución del optimizador, cobertura de automatización (proporción de agent) y la calificación de satisfacción del usuario.

\begin{table}[H]
\centering
\caption{Métricas para medir el ROI}
\begin{tabular}{p{0.22\textwidth}p{0.28\textwidth}p{0.28\textwidth}}
\toprule
Métrica & Ventana de observación & Dirección esperada \\
\midrule
Costo / consulta & Diaria & Reducir 10--20\% \\
Cobertura de automatización del agent & Semanal & Elevar a >70\% de la carga de trabajo \\
Satisfacción del usuario & Cada release & Elevar 0.2 puntos (sobre 5) \\
Éxito de invocación de herramientas & En tiempo real & > 95\% \\
\bottomrule
\end{tabular}
\end{table}

\subsection{Explicación del ROI}

El ROI no solo considera el ahorro en costos de la API, sino que también debe tomar en cuenta la reducción del tiempo del desarrollador. Mediante el análisis de trace de la Evidence Matrix es posible cuantificar el aumento en la velocidad de respuesta que `optimizer improvement` aporta al negocio y, a partir de ahí, calcular `value delivered / cost spent`.

\subsection{Resumen de la sección}

Las métricas de conversión en producto aseguran que la inversión en DSPy esté alineada con los objetivos del negocio; el cálculo del ROI requiere traducir las mejoras técnicas en cobertura del agent y en satisfacción del usuario.
\section{Agent Swarm y patrones de colaboración}

\subsection{Niveles de la estructura de Swarm}

El Swarm que menciona Omar generalmente incluye three tiers: 1) CLI/Coordinator, 2) execution agents, 3) verification agents. El Coordinator se encarga de la orchestration, los execution agents procesan las prompt/LLM calls, los verification agents revisan los outputs y asignan un `evidence rating`.

\begin{importantbox}{El valor de la separación de niveles}
La separación de niveles reduce la complejidad. Incluso si un execution agent comete un error en su output, el verification agent puede bloquear su paso por el pipeline en el nivel de la `Evidence Matrix`.
\end{importantbox}

\subsection{Colaboración de la cadena de herramientas}

Los tools del Swarm generalmente se exponen en forma de `MCP`; DSPy registra un `ToolDescriptor` al definir un Module. En runtime, el Coordinator invoca los tools en el orden que indica la `signature`, y registra los parámetros de la llamada y el valor de retorno en el evidence log.

\begin{itemize}
    \item `search\_tool`: invoca rápidamente el retrieval index
    \item `cohere\_tool`: invoca el LLM for consistent reasoning
    \item `validate\_tool`: run schema + rule-based checks sobre el tool output
\end{itemize}

\subsection{Resumen de la sección}

El éxito del Agent Swarm depende de la estructura jerárquica y la coordinación de herramientas: solo cuando el Coordinator, el Validator y el Executor cumplen cada uno su función es posible mantener bajo control el compound system.

\section{Observability Playbook}

\subsection{Arquitectura de Logging}

El stack de Logging se divide en tres capas: 1) request-level trace (prompt, module, timestamp), 2) evidence snapshot, 3) metric delta. Todos los three-tier log deben enviarse a un centralized observability system (como Grafana + Loki).

\subsection{Evidence-specific metrics}

Evidence metrics incluye prompt complexity, tool selection count, failure rate. Al comparar `evidence baseline` con `production run`, se puede descubrir cuándo comienza el optimizer drift. Se recomienda ejecutar `evidence replay` cada semana para verificar si las metrics están dentro del rango de +/-3%.

\subsection{Resumen de la sección}

Observability es el primer paso para volver operational a DSPy; solo si se puede ver el prompt/response/metric es posible realizar mejoras de bucle cerrado sobre el optimizer.

\section{Governance Playbook}

\subsection{Documentación y auditoría}

El sistema de gobernanza necesita incluir `DSPy.md`, Action Timeline, incident log. Contar con esta documentation y este audit trail garantiza la rendición de cuentas incluso en las actualizaciones del optimizer en sandbox.

\begin{itemize}
    \item `DSPy.md`: describe signature, owner, fallback
    \item `Action Timeline`: registra cada cambio del optimizer, trigger metric
    \item `Incident Log`: comparte la causa, los pasos de corrección, lessons learned
\end{itemize}

\subsection{Cumplimiento y aprobación}

Cada major update del optimizer necesita la aprobación del governance board, y los documentos de aprobación incluyen `Evidence Frame`, metric delta y el fallback plan. El board también establece una `review window` para el high-risk module (`tool` o `agent`).

\subsection{Resumen de la sección}

Governance Playbook convierte el DSPy panic en un proceso controlable, y protege al sistema de cambios excesivos mediante la cadena de documentación y aprobación.

\section{FAQ y síntesis de la experiencia}

\subsection{Preguntas frecuentes}

\begin{itemize}
    \item \textbf{P: ¿Para qué tamaño de proyecto es adecuado DSPy?} R: Es adecuado para todos los proyectos que necesitan combinar múltiples componentes de LLM, sobre todo cuando se quiere gobernar el prompt mediante código.
    \item \textbf{P: ¿Durante cuánto tiempo hay que conservar la Evidence Matrix?} R: Al menos 30 días o la ventana de tiempo del rollout anterior, para poder hacer regress reviews.
    \item \textbf{P: ¿Cómo se maneja el optimizer drift?} R: Primero se revisa la Evidence Matrix de manera retrospectiva, luego se ejecuta una prueba comparativa con el fallback prompt, y al mismo tiempo se registra el metric delta.
    \item \textbf{P: ¿Se necesita un DSL especializado?} R: DSPy ya incluye su propio DSL: basta con declarar los módulos y la signature en Python, sin necesidad de reinventar la rueda.
    \item \textbf{P: ¿Cómo se mide el optimizer improvements?} R: Se cuantifica mediante la combinación de tres indicadores principales: `automation coverage`, `latency` y `metric accuracy`.
\end{itemize}

\subsection{Lecciones extraídas de la experiencia}

La experiencia nos dice: 1) en cada release hay que entregar el evidence snapshot a QA; 2) el monitor necesita 24/7 alerting + weekly review; 3) la documentation no puede quedarse solo en la cabeza, debe escribirse como `DSPy.md` + `Incident Log`.

\subsection{Resumen de la sección}

El FAQ y la síntesis de la experiencia le dan al equipo respuestas rápidas a las preguntas frecuentes sobre el pipeline de DSPy, y evitan que se repitan los mismos errores en la práctica.

\section{Experimentos y validación}

\subsection{Principios de diseño experimental}

Cada DSPy experiment debe dejar en claro `hypothesis`, `metric` y `baseline`. Se recomienda un three-phase pipeline: 1) small-scale data → 2) evidence capture → 3) production dry-run. Compartir el evidence spreadsheet entre equipos permite que otros equipos repliquen el experimento.

\subsection{Simulation to Production}

Experimental validation no puede detenerse en offline benchmark, también se necesitan `sim-to-prod` steps. Cada optimizer change debe ejecutarse con 1000+ queries en el sandbox environment, y registrar `metric drift` y `cost per query`, para garantizar que el rollout sea controlable.

\begin{knowledgebox}{Simulation Pipeline}
El Simulation pipeline incluye `toy dataset` run, `evidence replay`, `sanity check`, y cada paso tiene un defined owner. Si el resultado de evidence replay difiere en más de 5\%, hay que volver a hacer tune o abort.
\end{knowledgebox}

\subsection{Resumen de la sección}

La sección de experimentos y validación ayuda a los equipos a hacer que los ajustes de DSPy sean reproducibles, y genera confianza antes del production rollout.

\section{Evidence Appendix}

\subsection{Ejemplos de Prompt y Metric}

\begin{table}[H]
\centering
\caption{Muestra de la Evidence Matrix}
\begin{tabular}{p{0.2\textwidth}p{0.35\textwidth}p{0.35\textwidth}}
\toprule
Tipo de prompt & Contenido generado & Métrica alcanzada \\
\midrule
Chain-of-thought question & Respuesta simulada de razonamiento en varios pasos & Accuracy 0.93 / cost 1.1x \\
Tool invocation with SQL & SQL generado automáticamente + Execution & Tool success 98\% \\
Reflective prompt & Resumen regenerado + calibration & BLEU 0.72 + Confidence 0.65 \\
Verifier prompt & Respuesta del modelo reconsiderada por el verifier & Drift drop 4\% after update \\
Fallback prompt & Ruta fallback de bajos recursos & Latency 0.8s / CLI coverage 43\% \\
Continuous context prompt & 6-shot progressive reasoning & Average tokens 450 \\
Action timeline prompt & Summarize pipeline decisions & Evidence rating 4.8/5 \\
\bottomrule
\end{tabular}
\end{table}

\subsection{Resumen de la sección}

Evidence Appendix muestra el registro correspondiente entre prompts y metrics, y es la fuente de primera mano para las decisiones del optimizer.
\section{Tabla de revisión de gobernanza}

\subsection{Eventos y acciones}

\begin{table}[H]
\centering
\caption{Tabla de revisión de gobernanza}
\begin{tabular}{p{0.18\textwidth}p{0.26\textwidth}p{0.26\textwidth}p{0.18\textwidth}}
\toprule
Evento & Descripción & Acción tomada & Owner \\
\midrule
Metric drift & Accuracy bajó de 0.94 a 0.87 & revertir la versión de optimizer, activar revisión & ML Ops \\
Tool failure & el SQL tool alcanzó el timeout 3 veces & aumentar el timeout & Tooling Team \\
Prompt change & Added fallback instructions & Update Action Timeline entry & Prompt Team \\
Governance review & Quarterly DSPy audit & Document evidence & Compliance \\
\bottomrule
\end{tabular}
\end{table}

\subsection{Resumen de la sección}

La tabla de revisión de gobernanza permite que el equipo, durante la revisión trimestral, entienda rápidamente los eventos, las acciones y los responsables, y así aporte context para la siguiente ronda de optimizer.

\section{Historia de recuperación de una falla}

\subsection{Contexto}

Después de un deployment, un User feedback spike indicó que el modelo se negaba a responder en cuestiones de cumplimiento; la Evidence Matrix señaló que `VerifyModule` arrojaba `score=0.55`.

\subsection{Pasos de recuperación}

\begin{enumerate}
    \item Revertir el optimizer result al previous stable commit
    \item Regenerar el fallback prompt, agregando la indicación ``if unsure, escalate to human''
    \item Registrar el rollback + reason con `Action Timeline`
    \item Trigger QA review to validate new prompt set
\end{enumerate}

\subsection{Resumen de la sección}

Esta recuperación de falla mostró cómo el evidence log, el Action Timeline y la documentación de governance se combinan para lograr una reversión rápida y una revisión posterior.

\section{Fragmento de código de ejemplo de DSPy}

\subsection{Pipeline declarativo}

\begin{lstlisting}[language=Python]
from dspy import Signature, Module, Optimizer, Evidence

class RetrieveModule(Module):
    signature = Signature("question", "context")
    def run(self, question):
        return search_knn(question, top_k=5)

class ReasonModule(Module):
    signature = Signature("context", "answer")
    def run(self, context):
        return llm_chat(f"Use the context to answer: {context}")

class VerifyModule(Module):
    signature = Signature("answer", "rating")
    def run(self, answer):
        score = verificator(answer)
        return {"score": score, "flags": []}

pipeline = [
    RetrieveModule(),
    ReasonModule(),
    VerifyModule()
]

optimizer = Optimizer(pipeline)
optimizer.search(metric="score", budget=100)

evidence = Evidence(pipeline, optimizer.result)
evidence.save("./evidence/latest.json")

# Additional helper routine used in experiments
def plan_and_execute(goal):
    plan = optimizer.plan(goal)
    for step in plan.steps:
        evidence.log(f"Executing {step.name}")
        result = step.module.run(step.input)
        evidence.attach(step.name, result)
        if result.get("score", 100) < 70:
            evidence.mark_failure(step.name)
            break
    return evidence.finalize()

# More DSL helpers for debugging
def replay_evidence(snapshot_path):
    snapshot = Evidence.load(snapshot_path)
    for entry in snapshot.entries:
        print(f"> {entry.prompt[:80]}... → {entry.metric}")
        if entry.metric < 0.6:
            print("   Needs review: ", entry.flags)
\end{lstlisting}

\subsection{Resumen de la sección}

Este fragmento de código muestra la estructura declarativa de DSPy combinada con optimizers, y ayuda al equipo a entender la secuencia real de llamadas detrás del compilador.

\section{Prompt Variation Grid}

\subsection{Tabla de variantes de Prompt}

\begin{table}[H]
\centering
\caption{Prompt Variation Grid}
\begin{tabular}{p{0.2\textwidth}p{0.3\textwidth}p{0.3\textwidth}}
\toprule
Variante & Descripción & metric objetivo \\
\midrule
Bare prompt & pregunta directa & Baseline accuracy \\
Chain-of-thought & discusión de varios pasos & Accuracy + explanation score \\
ReAct & Prompt + tool hints & Tool success rate \\
Verifier double-check & format response → verify & False positive drop \\
Few-shot 3 examples & include 3 QA pairs & Consistency \\
Hard-coded instructions & specify style + tone & Style adherence \\
Responder hint & tell assistant to output bullet list & Structure quality \\
Planner injection & ask agent to plan first & Planning coverage \\
Tool-specific prompt & instruct tool usage order & Tool selection accuracy \\
Fallback fallback & add fallback path for rejects & Robust failure handling \\
Calibration prompt & ask agent to self-score & Confidence calibration \\
Meta evidence prompt & log prompt/res metric & Observability lookups \\
\bottomrule
\end{tabular}
\end{table}

\subsection{Resumen de la sección}

Prompt Variation Grid ayuda al equipo a cuantificar el desempeño de las distintas estrategias, y le da al optimizer más candidatos en lugar de un único prompt.

\section{Estudio de caso de despliegue}

\subsection{Contexto del caso}

Cierto cliente financiero necesitaba un multi-step reasoning agent para procesar investigaciones de cumplimiento. El equipo adoptó DSPy y encadenó la recuperación, la inferencia del LLM, la SQL tool y el Verify module. El principal riesgo antes del despliegue eran las vulnerabilidades de tool y el metric drift.

\subsection{Medidas de control}

\begin{itemize}
    \item Cada día se guarda el Evidence Snapshot en el central repo y se registra la optimizer version mediante `git tag`
    \item Configurar `Action Timeline` para forzar la contención del `fallback path` antes del production release
    \item Monitorear el tool success rate y activar automáticamente la secuencia `tool redrive` después de cada falla
    \item Evaluación multimodelo: ejecutar simultáneamente two DSPy variants en una prueba A/B y llevar la better variant a production
\end{itemize}

\subsection{Resumen de la sección}

Este caso muestra cómo el DSPy pipeline se implementa en escenarios complejos: mediante evidence logging, action timeline y dual evaluation se mantiene el riesgo bajo control.

\section{Glosario \& cadena de herramientas}

\begin{description}
    \item[Signature] Declaración de entrada/salida del módulo; DSPy depende de ella para construir pipelines componibles.
    \item[Evidence Matrix] Tabla de registro de prompt/response/metric.
    \item[Optimizer] Stack que busca automáticamente prompts/ejemplos.
    \item[Action Timeline] Registro de los ajustes del pipeline y su owner.
    \item[MCP] Modular Control Protocol, usado para el handshake entre agents.
    \item[Tool Descriptor] Metadata que describe los parámetros que acepta un tool.
    \item[Fallback] Prompt/pipeline de contingencia que el monitoreo activa.
    \item[Metric Drift] Descenso de las metrics en production, que normalmente requiere retrain/rollback.
    \item[Governance Log] Documento que registra incidentes, responsables y conclusiones de la retrospectiva.
    \item[Evidence Snapshot] Captura en tiempo de ejecución de prompt/response/carbon copy, usada para su reproducción posterior.
\end{description}

\subsection{Resumen de la sección}

El glosario y la explicación de la cadena de herramientas le permiten al equipo mantener un lenguaje consistente en las reuniones interfuncionales, y evitar el riesgo de que ``cada quien entienda algo distinto''.
\section{Casos prácticos y ruta de implementación}

\subsection{Caso: sistema Multi-hop DSPy}

Un pipeline multi-hop típico primero recupera, después razona, luego verifica y después invoca herramientas. DSPy te permite abstraer cada uno de estos pasos por separado en un module, y luego usar un optimizer para ayudarte a buscar combinaciones de prompt y ejemplos. En la práctica te encontrarás con problemas de sincronización en la línea de tiempo entre `chain-of-thought` y `tool invocation`.

\subsection{Gobernanza y documentación}

Se recomienda agregar un documento `DSPy.md` nuevo en el repo, que incluya la versión de deployment, los enlaces del Action Timeline y del Evidence Matrix, así como la security checklist. En cada release, `DSPy.md` se incluye como parte de la release note y se entrega a QA para su revisión.

\begin{importantbox}{La documentación es gobernanza}
El código de DSPy tiene dos capas más que un pipeline normal: DSL + optimizer log. Si no se documenta, aparecerá entre los desarrolladores la fricción de ``no sé por qué el optimizer cambió el prompt''. La documentación debe incluir: signature, metric, condiciones de fallback, contacto del owner.
\end{importantbox}

\subsection{Resumen de la sección}

Descomponer los sistemas compuestos de gran escala en casos reutilizables, y mantener para ellos documentación en texto plano, es un paso clave para pasar de la prueba de concepto a la operación real.

\section{Direcciones de investigación abiertas}

\begin{itemize}
    \item Optimizadores más potentes: ¿cómo optimizar eficientemente en espacios de búsqueda más grandes?
    \item Optimización conjunta entre módulos: ¿cómo optimizar simultáneamente varios módulos dentro del pipeline?
    \item Integración con arquitecturas de Agent: ¿cómo aplicar el paradigma de optimización de DSPy a sistemas de Agent como ReAct?
    \item Ecosistema abierto: fomentar que la comunidad aporte nuevos módulos, optimizadores y métricas de evaluación
\end{itemize}

\section{Síntesis y ampliación}

\subsection{Puntos clave}

\begin{enumerate}
    \item Los sistemas de IA compuestos superan los límites de capacidad de una sola llamada al LLM al combinar varios componentes
    \item El prompting manual es frágil, no componible y difícil de optimizar en sistemas compuestos
    \item DSPy ofrece un paradigma de programación declarativa: definir la estructura + compilar el prompt automáticamente
    \item Signature, Module, Optimizer y Metric son las cuatro abstracciones centrales
    \item Analogía con un compilador: lenguaje de alto nivel (programa DSPy) → implementación de bajo nivel (prompt optimizado)
\end{enumerate}

\subsection{Lecturas adicionales}

\begin{itemize}
    \item Khattab et al., ``DSPy: Compiling Declarative Language Model Calls into Self-Improving Pipelines,'' ICLR 2024.
    \item Khattab et al., ``Demonstrate-Search-Predict: Composing Retrieval and Language Models for Knowledge-Intensive NLP,'' 2023.
\end{itemize}

\subsection{Tabla resumen global}

\begin{table}[H]
\centering
\begin{tabular}{p{0.22\textwidth}p{0.25\textwidth}p{0.25\textwidth}}
\toprule
Tema & Mejora de capacidad & Riesgo de ingeniería \\
\midrule
Compound Systems & Múltiples módulos en pipeline & Las vulnerabilidades de composición se vuelven fallas en cadena \\
Compilación de DSPy & Generación declarativa de prompts & Espacio de búsqueda grande, se necesitan resultados rastreables \\
Evidence/Monitoring & Permite observar cada llamada & La ausencia de monitoreo amplifica la degradación \\
\bottomrule
\end{tabular}
\caption{Tabla de valor y riesgo de la Lecture 08}
\end{table}

\subsection{Acciones futuras}

\begin{itemize}
    \item Establecer un baseline y un rollout log para cada versión del optimizer
    \item Vincular la Evidence Matrix con el sistema de alertas para detectar automáticamente el prompt drift
    \item Mantener actualizado el Action Timeline para ubicar rápidamente a los responsables
\end{itemize}


\end{document}
